\section{Introduction}
\label{sec:intro}
Battery-powered three-wheelers (easy-bikes, Mishuks, e-rickshaws) have become a mass transport mode in Bangladesh. A CLEAN/BWGED study estimates about 1.4 million easy-bikes consuming roughly 4{,}905 GWh of electricity per year \cite{clean2024}, whereas Pure Earth puts the number of electric three-wheelers at 3--4 million \cite{pureearth_smep}; we treat both as scenarios. Lead-acid batteries dominate \cite{mordor2026} and are short-lived: a battery set costing more than BDT 72{,}000 is reported to last 8--11 months \cite{fe2024}, and a 2026 field study of 150 drivers in Gopalganj and Khulna documents prolonged hub charging, frequent overloading and lead-acid degradation \cite{chowdhury2026}.

This short life has three consequences. First, replacement is a large part of operator cost. Second, internal resistance and charger losses waste electricity, and each replacement carries manufacturing emissions. Third, used lead-acid batteries (ULABs) are broken and smelted largely by the informal sector, a primary contributor to lead-pollution hotspots \cite{esdo2021}. The national public-health bulletin of IEDCR reports that about 36 million Bangladeshi children (60\%) are poisoned by lead and, citing a World Bank study in The Lancet, that lead causes a loss of about 20 million IQ points and 140{,}000 cardiovascular deaths in adults; the 2019 economic burden was US\$28.6 billion, or 6--9\% of GDP \cite{das2026}. Pure Earth reports about 125 kg of lead in one easy-bike's batteries and about 167{,}000 t of lead waste per year from recycling them \cite{fe2024,pureearth_smep}.

LFP batteries offer higher round-trip efficiency, longer cycle life and no lead, but cost more. Laboratory cycle-life figures suggest they should be cheaper over a bike's life; whether drivers' actual prices and replacement intervals support this is an empirical question. Bangladesh is also hot: monthly mean 2\,m temperature in the study area ranges from \Tmin\ to \Tmax\,$^\circ$C (\figref{temp}), and a 42\,$^\circ$C day has been reported at Chuadanga \cite{heat2025}. Heat accelerates ageing of every chemistry, but not equally.

\textbf{Research gap.} Existing work compares lead-acid and lithium-ion for stationary storage \cite{keshan2016,yudhistira2022,jasper2025}, studies cell ageing \cite{yagci2026}, thermal safety \cite{schoberl2024}, charging infrastructure \cite{hasan2020}, or easy-bike operation and energy use for lead-acid only \cite{chowdhury2026,peakdemand}. In web searches of publisher and data sites in October 2026 we found no study that anchors a temperature-driven chemistry comparison for Bangladesh easy-bikes on drivers' own purchase prices, replacement intervals and resale values, and that adds carbon and lead-flow accounting. \todo{[State the databases, search terms and dates of your final literature search.]}

\textbf{Contributions.} (i) A field survey of \nSurvey{} drivers (October 2026) whose cleaned data supply the cost, life, resale and distance inputs. (ii) A reproducible MATLAB/Simulink model, driven by NASA POWER temperature, with analytic and block-diagram versions cross-checked. (iii) A calibration and validation design that exposes how far laboratory cycle lives are from field practice (\figref{validlife}). (iv) Break-even price and life for LFP, sensitivity, carbon and lead-flow accounting, and the operator support implied by conversion scenarios.

Section~\ref{sec:data} reviews the literature and inputs, Section~\ref{sec:method} describes the survey and model, Section~\ref{sec:results} gives results, Section~\ref{sec:disc} discusses implications and limits, and Section~\ref{sec:concl} concludes.
